% Generated draft: requires mathematical and source review.
\documentclass[11pt]{article}
\usepackage[margin=25mm]{geometry}
\usepackage{fontspec}
\setmainfont{Latin Modern Roman}
\usepackage{amsmath,amssymb,mathtools,unicode-math}
\setmathfont{Latin Modern Math}
\usepackage{enumitem,needspace,xurl,hyperref}
\setlength{\parindent}{0pt}
\setlength{\parskip}{5pt}
\setlength{\emergencystretch}{3em}
\newcommand{\N}{\mathbb{N}}
\newcommand{\Z}{\mathbb{Z}}
\newcommand{\Q}{\mathbb{Q}}
\newcommand{\R}{\mathbb{R}}
\newcommand{\C}{\mathbb{C}}
\newcommand{\F}{\mathbb{F}}
\newcommand{\A}{\mathbb{A}}
\newcommand{\PP}{\mathbb{P}}
\newcommand{\E}{\mathbb{E}}
\newcommand{\eps}{\varepsilon}
\newcommand{\dd}{\,\mathrm{d}}
\begin{document}
\section*{UnsolvedMath ID: 10400139}
\textbf{Source problem code:} catalogue entry
\subsection{Definitions and mathematical statement}
Let \(\mathcal L\) denote the class of links in \(S^3\) to which the source invariants apply, with their required equivalence and auxiliary-data conventions. For each allowed integer parameter \(N\), put
\[
 \omega_N=\exp(2\pi i/N),\qquad
 Q_N(L)=K_N(S^3,L),\qquad C_N(L)=J_N(L;\omega_N),
\]
where \(J_N\) is the coloured Jones polynomial in the source's convention, normalized by \(J_N(U)=1\) for the unknot \(U\). The problem, attributed to S. Baseilhac and R. Benedetti, is to determine the precise relationship between the functions \(Q_N\) and \(C_N\) of links. In particular, the comparison must specify its range of \(N\) and links and any normalization or auxiliary choices on which it depends. The source does not assert a particular equality between these functions. Its excerpt does not define \(K_N\), the allowed parameters, or the orientation, framing and phase conventions, so no unique comparison identity can be extracted from it.
\subsection{Short English statement}
Determine exactly how the quantum hyperbolic invariant \(K_N(S^3,L)\) is related, as a function of links, to the coloured Jones polynomial evaluated at \(\exp(2\pi i/N)\) and normalized to be \(1\) on the unknot.
\subsection{Sources}
\begin{itemize}
\item S. Baseilhac and R. Benedetti, attribution in Problem 7.24 of T. Ohtsuki, \emph{Problems on invariants of knots and 3-manifolds} (2002), supplied PDF page 115. \url{https://msp.org/gtm/2002/04/gtm-2002-04-024s.pdf}. Recorded access: 2026-07-29; extraction: pdf-font-structure.
\item Local UnsolvedMath registry, \texttt{data/unsolvedmath/problems.json}, permanent ID \(10400139\), exact source code \texttt{catalogue entry}. The supplied research summary identifies the QHI/coloured Jones comparison and labels it OPEN-TRIAGE; it supplies no resolving reference.
\end{itemize}
\end{document}
